\documentclass[11pt,fleqn]{article} 
\usepackage[margin=0.8in, head=1in]{geometry} 
\usepackage{amsmath, amssymb, amsthm}
\usepackage{fancyhdr} 
\usepackage{palatino, url, multicol}
\usepackage{graphicx} 
\usepackage[all]{xy}
\usepackage{polynom} 
\usepackage{pdfsync}
\usepackage{enumerate}
\usepackage{framed}
\usepackage{setspace}
\usepackage{array,tikz}
\pagestyle{fancy} 
\lhead{Linear Algebra}
\chead{Some Practice Problems }
%\rhead{\Large{Name: \underline{\hspace{2in}}}}

\newcommand{\vtr}[1]{\textbf{#1}}

\def\vectwo#1#2{\begin{bmatrix}#1\\#2\end{bmatrix}}
\def\vecthree#1#2#3{\begin{bmatrix}#1\\#2\\#3\end{bmatrix}}
\def\vecfour#1#2#3#4{\begin{bmatrix}#1\\#2\\#3\\#4\end{bmatrix}}
\def\bbm{\begin{bmatrix}}
\def\ebm{\end{bmatrix}}

% macros
\usepackage{amssymb}
\newcommand{\bA}{\mathbf{A}}
\newcommand{\bB}{\mathbf{B}}
\newcommand{\bE}{\mathbf{E}}
\newcommand{\bF}{\mathbf{F}}
\newcommand{\bJ}{\mathbf{J}}

\newcommand{\bb}{\mathbf{b}}
\newcommand{\bc}{\mathbf{c}}
\newcommand{\br}{\mathbf{r}}
\newcommand{\bv}{\mathbf{v}}
\newcommand{\bw}{\mathbf{w}}
\newcommand{\bx}{\mathbf{x}}

\begin{document}
\renewcommand{\headrulewidth}{0pt}
\newcommand{\blank}[1]{\rule{#1}{0.75pt}}
\renewcommand{\d}{\displaystyle}

%\vspace*{-0.7in}

\begin{enumerate}
\item
Consider the vectors

$$\mathbf a=(1,-1,3),\ \ \mathbf b=(2,1,1)$$
\begin{enumerate}

\item Are $\mathbf a$ and $\mathbf b$ orthogonal? 
\vskip .7in

\item Give a third vector $\mathbf c$ that is orthogonal to $\mathbf a$ and one unit long. 
\vskip 1.in

\item How many correct answers are there to part (b)?
\vskip .7in

\end{enumerate}

\item  Consider the equation
 $$\begin{bmatrix} 1 &-2\\ 3& 2\end{bmatrix}\begin{bmatrix} x\\y\end{bmatrix}=\begin{bmatrix} 3\\-4\end{bmatrix}.$$
\begin{enumerate}
\item Sketch (approximately) the column picture for this equation. 
\vskip 1.5in
\item In a sentence or two, explain how this picture indicates this equation is solvable.

\vskip .8in
\item Change only one entry in the matrix so that the equation cannot be solved.
\vskip 1in
\end{enumerate}
 
 \item Solve the following, using elimination on the augmented matrix, and back substitution:
$$\begin{pmatrix} 1&1&1\\-1&1&-2\\0&-6&1\end{pmatrix}\mathbf x=\begin{pmatrix}2\\1\\-7\end{pmatrix}.$$

\vskip 2.5in

\item 
\begin{enumerate}
\item Give the two $3\times3$ elimination matrices that put the above system in upper-triangular form.
\vskip 1.in

\item Give \emph{one} matrix that does both these elimination steps at once. (Make sure you think about the correct order of the steps.)
\end{enumerate}

\vskip 3.in


\item What should $z$ be to make the the  (2,3)-entry of  following  matrix product be $0$?
$$\begin{pmatrix} 1&2&z\\2&z&3\\z&3&4\end{pmatrix}
\begin{pmatrix} 1&1&0\\0&2&1\\-1&1&-1\end{pmatrix}
$$
 

\vskip 1.2in


\item Find $A^{-1}$, or show it doesn't exist, for
$$
A=\begin{pmatrix} 1&0&1\\2&0&1\\-1&2&-1\end{pmatrix}.
$$
 

\vskip 3.in


\item   Give an $LU$ factorization of 
$$A=\begin{pmatrix}2&1&-2\\-4&-3&5\\0&-1&4\end{pmatrix}.$$



\vskip 2.5in


\item Suppose the $LU$ factorization of a matrix $A$ is 
$$A=\begin{pmatrix}1 &0&0\\-2 &1&0\\5&-1&1\end{pmatrix}\begin{pmatrix}1& 0 &-1\\0 &-2&1\\0&0&3\end{pmatrix}.$$
Use this to solve $A\mathbf x =\mathbf b$ for $\mathbf b=(-1,6,-3)$ by solving two triangular systems. Do NOT compute $A$ to do this problem!

\vskip 1.in


%%%%%%%%%%%%%
%%% END RHODES
%%% BEGIN BUELER
%%%%%%%%%%%%%

\item  Here are three equations in two unknowns:
\begin{align*}
2 x + 2 y &= 6 \\
x - 3 y &= -1 \\
4 x + y &= 0
\end{align*}
We can also write this as $A \bx = \bb$ where $A$ is a $3\times 2$ matrix and $\bx = \begin{bmatrix} x \\ y \end{bmatrix}$, $\bb = \begin{bmatrix} 6 \\ -1 \\ 0 \end{bmatrix}$ are vectors.

	\begin{enumerate}
	\item Sketch the ``row picture'': sketch each equation as a line in the $(x,y)$ plane.  Do they intersect?

	\vfill
	\item Sketch the ``column picture'': sketch each column of $A$, and also $\bb$, in three-dimensional space.  Will you be able to find a linear combination of the columns of $A$ which gives $\bb$?

\vfill


	\item  Change one entry of the right side so that the linear system does have a solution, and find that solution.

\vspace{2.5in}
	\end{enumerate}
\item
	\begin{enumerate}
	 \item  Consider the new linear system $A\bx=\bb$ where
    $$A = \begin{bmatrix} 1 & 0 & -1 \\ 2 & 2 & 0 \\ 3 & 2 & 1 \end{bmatrix}, \quad \bx = \begin{bmatrix} 0 \\ 4 \\ 6 \end{bmatrix},$$
and $\bx = (x_1,x_2,x_3)$ is unknown (\emph{for now}).  Sketch the row picture, that is, sketch each of the three equations as a plane.  Note it is easier to sketch each plane on separate axes; show three sketches.
\vfill

	\item  What is the solution of the system in part \textbf{(i)}?
\vspace{1.0in}
	\end{enumerate}


\item  Consider the following linear system $A\bx = \bb$:
\begin{align*}
2 x_1 + x_2 - 9 x_3 &= -6 \\
4 x_1 - 3 x_2 -21 x_3 &= -20 \\
-6 x_1 - 13 x_2 + 24 x_3 &= 5
\end{align*}

Solve the linear system by \emph{elimination} and then \emph{back-substitution}.  Use the standard algorithm.  Show your work, and in particular show, as an intermediate stage, the triangular system which you get after elimination.
\vfill


\vspace{3in}
\item  From problem above, 
	\begin{enumerate} 
	\item what elimination matrices $E_{21},E_{31}E_{32}$ did the row operations?
	\vfill

	\item  What are the inverses $L_{21}=E_{21}^{-1}$, $L_{31}=E_{31}^{-1}$, $L_{32}=E_{32}^{-1}$?
	\vfill

	\item  What numbers were the pivots?  What is the determinant of A?
	\vfill

	\item  The computation can regarded as factoring $A=LU$.  What lower triangular matrix $L$ and upper triangular matrix $U$ were computed?
\vfill

	\item   Multiply $LU$ and confirm you get the original matrix $A$.
\vspace{0.5in}
	\end{enumerate}


\end{enumerate}
\end{document}